\chapter{Running a Multi-Strategy Book}\label{ch:s1:running-a-multi-strategy-book}

In the first week of August 2007 many quantitative equity funds lost at once. Khandani and Lo, rebuilding the episode from factor and transaction data,
found a market-wide deleveraging: two unwinds, on 1 August and from the open of 6 August, beginning with financials, long book-to-market and short
earnings momentum, the positions many funds held under different names. A firm that runs many strategies is exposed to the same thing inside its own
walls. On this chapter's synthetic firm of ten \omterm{def:s1:running-a-multi-strategy-book:multi}{pods}, five share hidden factors that crash together three times in twenty years; detecting the overlap
and capping it cuts the firm's largest drawdown from 43.5\% to 35.7\% at 10\% volatility, and a 5\% \omterm{def:s1:running-a-multi-strategy-book:limit}{drawdown limit} on each \omterm{def:s1:running-a-multi-strategy-book:multi}{pod} stops 91.6\% of healthy
\omterm{def:s1:running-a-multi-strategy-book:multi}{pods} within a year. The build is \texttt{firm.multistrat}.

\section{Allocation across strategies}

\begin{definition}[Multi-strategy book, pod]\label{def:s1:running-a-multi-strategy-book:multi}
A \emph{multi-strategy book}\index{multi-strategy book} is a firm's combined portfolio of several strategies run by separate teams, allocated capital or
risk by a central function that nets their trades and controls their combined risk. A \emph{pod}\index{pod} is one such team with its own capital, risk
limits and P\&L.
\end{definition}

\begin{definition}[Strategy correlation]\label{def:s1:running-a-multi-strategy-book:correlation}
\emph{Strategy correlation}\index{strategy correlation} is the correlation between the P\&L streams of two strategies; measured over normal periods it can
be low while the strategies share an exposure that dominates in stress, so a firm estimates it both from P\&L and from the positions' factor exposures.
\end{definition}

The firm's first decision is how much risk each \omterm{def:s1:running-a-multi-strategy-book:multi}{pod} gets. Book~7, chapter~26's risk budgeting gives the default: weights inversely proportional to each
\omterm{def:s1:running-a-multi-strategy-book:multi}{pod}'s trailing volatility, so that each contributes similar risk, rebalanced monthly (\cref{lst:s1:running-a-multi-strategy-book:firm}). The synthetic
firm has ten \omterm{def:s1:running-a-multi-strategy-book:multi}{pods} at 10\% volatility each, with Sharpe ratios from 0.5 to 1.0. \omterm{def:s1:running-a-multi-strategy-book:multi}{Pods} 1 to 3 load 0.4 on a shared factor, \omterm{def:s1:running-a-multi-strategy-book:multi}{pods} 4 and 5 on another; the
factors are as quiet as any other noise except in three ten-day crashes, each a fall of six \omterm{def:s1:running-a-multi-strategy-book:multi}{pod} volatilities. \omterm{def:s1:running-a-multi-strategy-book:multi}{Pod} 10's edge dies after ten years.

\section{Netting and shared risk}

Measured on twenty years of daily P\&L, the three \omterm{def:s1:running-a-multi-strategy-book:multi}{pods} on the first factor have an average correlation of 0.22; a \omterm{def:s1:running-a-multi-strategy-book:multi}{pod} on neither factor, $-0.004$
with the first. The correlation understates what they share. In each crash the three \omterm{def:s1:running-a-multi-strategy-book:multi}{pods} lose 2.4 to 2.6 of their annual volatilities together, 24\% to
26\% of capital, in ten days. The same trade under different names is invisible in correlations measured over calm years and obvious in the \omterm{def:s1:running-a-multi-strategy-book:multi}{pods}'
positions: a risk model (Book~7, chapter~24) maps each \omterm{def:s1:running-a-multi-strategy-book:multi}{pod}'s holdings to the factors, and the firm can cap the total exposure to each.

\begin{center}
\small
\begin{tabular}{@{}lccc@{}}
\toprule
firm at 10\% volatility & no overlap control & factor cap 1.0 & factor cap 0.5\\
\midrule
Sharpe ratio & 1.45 & 1.52 & 1.70\\
largest drawdown & 43.5\% & 42.8\% & 35.7\%\\
loss in the three crashes & 31.5, 31.7, 33.6\% & 30.1, 30.1, 32.3\% & 21.6, 20.7, 22.9\%\\
\bottomrule
\end{tabular}
\end{center}

The cap limits each shared factor's total exposure to the equivalent of one (or half) equal-weight \omterm{def:s1:running-a-multi-strategy-book:multi}{pod}, scaling down the \omterm{def:s1:running-a-multi-strategy-book:multi}{pods} that carry it. Capping at
half a \omterm{def:s1:running-a-multi-strategy-book:multi}{pod} cuts each crash by about a third and raises the firm's Sharpe ratio from 1.45 to 1.70 (\cref{fig:s1:running-a-multi-strategy-book:paths}),
because the crashes were the main cost of the overlapping \omterm{def:s1:running-a-multi-strategy-book:multi}{pods}. The losses are large because a firm of ten nearly uncorrelated \omterm{def:s1:running-a-multi-strategy-book:multi}{pods} must lever them to
reach 10\% volatility, and leverage multiplies whatever they share.

\begin{omfigure}
\begin{tikzpicture}
\begin{axis}[width=11cm,height=5.2cm,xlabel={\small year},ylabel={\small cumulative return (\%)},tick label style={font=\footnotesize},
  xmin=1,xmax=20,ymin=-20,ymax=340,grid=major,grid style={gray!25},legend pos=north west,legend style={font=\scriptsize},table/col sep=comma]
\fill[omExo!15] (axis cs:5.952,-20) rectangle (axis cs:5.992,340);
\fill[omExo!15] (axis cs:14.286,-20) rectangle (axis cs:14.325,340);
\fill[omExo!15] (axis cs:17.460,-20) rectangle (axis cs:17.500,340);
\addplot[omDef,thick] table[x=year,y=none]{figdata/strategies-1/28-running-a-multi-strategy-book/paths.csv};
\addplot[omThm,thick] table[x=year,y=cap]{figdata/strategies-1/28-running-a-multi-strategy-book/paths.csv};
\legend{no overlap control,shared factors capped at half a pod}
\end{axis}
\end{tikzpicture}

\omcaption{The synthetic firm of ten \omterm{def:s1:running-a-multi-strategy-book:multi}{pods}, scaled to 10\% volatility, with and without a cap on the \omterm{def:s1:running-a-multi-strategy-book:multi}{pods}' shared factor exposures; cumulative sums of
daily returns over twenty years, the shaded bands marking the three crashes of the shared factors. Data: \texttt{s1\_multistrat.paths}.}\label{fig:s1:running-a-multi-strategy-book:paths}
\end{omfigure}

The same central view saves costs. When \omterm{def:s1:running-a-multi-strategy-book:multi}{pods} trade the same instruments, the firm executes only the net: Book~7, chapter~27's internal crossing. Ten
\omterm{def:s1:running-a-multi-strategy-book:multi}{pods} trading 200 instruments independently, each trading a fifth of them, have 34.6\% of their gross volume cancel inside the firm; if all trade all,
68.5\%, close to $1 - 1/\sqrt{10}$.

\section{Turning strategies off}

\begin{definition}[Drawdown limit]\label{def:s1:running-a-multi-strategy-book:limit}
A \emph{drawdown limit}\index{drawdown limit} stops or cuts a strategy's risk when its loss from its previous peak exceeds a threshold, a share of its
allocated capital; it trades the cost of stopping good strategies in bad luck against the cost of running dead ones too long.
\end{definition}

Every stop rule makes two errors. It stops \omterm{def:s1:running-a-multi-strategy-book:multi}{pods} whose edge is alive, and it lets dead ones run. \Cref{fig:s1:running-a-multi-strategy-book:stops} shows
the first: a \omterm{def:s1:running-a-multi-strategy-book:multi}{pod} with a Sharpe ratio of 0.7 at 10\% volatility hits a 5\% drawdown within a year on 91.6\% of 20\,000 simulated paths, a 10\% drawdown on
34.5\%, 20\% on 2.1\%. A limit is measured in the \omterm{def:s1:running-a-multi-strategy-book:multi}{pod}'s own volatilities: 5\% is half a year's standard deviation, well inside ordinary noise. In the
synthetic firm, where stopped \omterm{def:s1:running-a-multi-strategy-book:multi}{pods} restart each year, a 5\% limit stops 85.8 live \omterm{def:s1:running-a-multi-strategy-book:multi}{pods} for every hundred pod-years; at 20\% it still stops 14.7,
because the shared crashes push several \omterm{def:s1:running-a-multi-strategy-book:multi}{pods} through it together.

\begin{omfigure}
\begin{tikzpicture}
\begin{axis}[width=10cm,height=5cm,xlabel={\small drawdown limit (\% of the pod's capital)},ylabel={\small \% of live pods stopped in a year},
  tick label style={font=\footnotesize},xmin=3,xmax=32,ymin=0,ymax=100,grid=major,grid style={gray!25},legend pos=north east,
  legend style={font=\scriptsize},table/col sep=comma]
\addplot[omThm,thick,mark=*] table[x=limit_pct,y=paths_pct]{figdata/strategies-1/28-running-a-multi-strategy-book/stops.csv};
\addplot[omDef,thick,mark=square*] table[x=limit_pct,y=firm_pct]{figdata/strategies-1/28-running-a-multi-strategy-book/stops.csv};
\legend{one pod (Sharpe 0.7) on 20\,000 paths,pods of the synthetic firm}
\end{axis}
\end{tikzpicture}

\omcaption{False stops: the share of \omterm{def:s1:running-a-multi-strategy-book:multi}{pods} with a live edge that hit a \omterm{def:s1:running-a-multi-strategy-book:limit}{drawdown limit} within a year, for a single \omterm{def:s1:running-a-multi-strategy-book:multi}{pod} with a Sharpe ratio of 0.7 at 10\%
volatility (simulated paths) and for the synthetic firm's \omterm{def:s1:running-a-multi-strategy-book:multi}{pods}, whose shared crashes add stops. Data: \texttt{s1\_multistrat.limits}.}\label{fig:s1:running-a-multi-strategy-book:stops}
\end{omfigure}

The second error is the dead \omterm{def:s1:running-a-multi-strategy-book:multi}{pod}. Its Sharpe ratio falls from 0.9 to $-0.3$ after ten years; the 5\% limit stops it 98 trading days later, the 10\% limit
after 307, the 15\% limit after 979, and wider limits never, in the ten years left. A tight limit catches the dead \omterm{def:s1:running-a-multi-strategy-book:multi}{pod} quickly by catching nearly every
\omterm{def:s1:running-a-multi-strategy-book:multi}{pod} quickly. The firm's Sharpe ratio across limits (1.49 at 5\%, 1.74 at 10\%, 1.61 at 15\%, 1.48 at 20\%) peaks in the middle, and in a sample this size
that peak is not a robust optimum.

\section{Capital, drawdown limits and the pod model}

The \omterm{def:s1:running-a-multi-strategy-book:limit}{drawdown limit} has a theory. Grossman and Zhou solved the problem of an investor who never wants wealth to fall below a fixed fraction of its
running maximum: the optimal risky position is proportional to the distance to that floor. Risk falls as losses approach the floor instead of switching
off at it, and it rises again as the surplus is rebuilt. A \omterm{def:s1:running-a-multi-strategy-book:multi}{pod} model that halves a \omterm{def:s1:running-a-multi-strategy-book:multi}{pod}'s risk at one threshold and stops it at a second is a two-step
version of that rule. The central lessons of the chapter follow. Limits belong in units of each \omterm{def:s1:running-a-multi-strategy-book:multi}{pod}'s volatility and expected Sharpe ratio, not in
round percentages. Shared exposures must be measured from positions, not from P\&L correlations. And the firm's capital is safest when drawdowns are
controlled continuously, as Grossman and Zhou's rule does, rather than by a cliff that stops many good \omterm{def:s1:running-a-multi-strategy-book:multi}{pods} to catch a few bad ones.

\section{Strategy files}

\begin{strategyfile}{Multi-strategy allocation with overlap control}\label{strat:s1:running-a-multi-strategy-book:overlap}
\sfield{Who pays you, and why} The \omterm{def:s1:running-a-multi-strategy-book:multi}{pods}' combined edges, with diversification when their risks are truly different.
\sfield{Instruments and venues} Whatever the \omterm{def:s1:running-a-multi-strategy-book:multi}{pods} trade.
\sfield{Signal} \omterm{def:s1:running-a-multi-strategy-book:multi}{Pods}' trailing volatility and their positions' factor exposures.
\sfield{Sizing and execution} Risk budgets by inverse volatility; caps on each shared factor; trades netted across \omterm{def:s1:running-a-multi-strategy-book:multi}{pods}.
\sfield{Costs} Lower through netting; a risk model and a central execution desk.
\sfield{How it dies} Crowded factors shared by many firms, as in August 2007.
\sfield{Horizon, capacity, infrastructure} Continuous; position-level data from every \omterm{def:s1:running-a-multi-strategy-book:multi}{pod}.
\sfield{Backtest honestly} Exposures from positions at each date; crash episodes in the sample.
\sfield{Sources} Khandani and Lo (2011); this chapter: drawdown 43.5\% to 35.7\% with a cap of half a \omterm{def:s1:running-a-multi-strategy-book:multi}{pod}.
\end{strategyfile}

\begin{strategyfile}{Drawdown-controlled exposure}\label{strat:s1:running-a-multi-strategy-book:drawdown}
\sfield{Who pays you, and why} Nobody: it is a risk rule that trades expected return for a floor on losses.
\sfield{Instruments and venues} Any strategy or the firm.
\sfield{Signal} The distance from the running peak to a floor.
\sfield{Sizing and execution} Risk proportional to the surplus over the floor, or cut in steps.
\sfield{Costs} False stops and the return given up after them.
\sfield{How it dies} Limits set in percentages rather than volatilities; shared crashes that hit many \omterm{def:s1:running-a-multi-strategy-book:multi}{pods} together.
\sfield{Horizon, capacity, infrastructure} Continuous; daily P\&L.
\sfield{Backtest honestly} Count false stops as well as avoided losses.
\sfield{Sources} Grossman and Zhou (1993); this chapter: 91.6\% of healthy \omterm{def:s1:running-a-multi-strategy-book:multi}{pods} stopped in a year at a 5\% limit.
\end{strategyfile}

\section{Tutorial: five at once}\label{tut:s1:running-a-multi-strategy-book:tut}

\begin{tutorial}
\textbf{Goal.} Run a firm of ten \omterm{def:s1:running-a-multi-strategy-book:multi}{pods} with hidden shared factors, allocate with and without overlap control, apply \omterm{def:s1:running-a-multi-strategy-book:limit}{drawdown limits} and measure false
stops and the dead \omterm{def:s1:running-a-multi-strategy-book:multi}{pod}'s survival, and net trades. \textbf{End state:} the tables and the two figures.
\begin{tutsteps}
\item \textbf{Allocation and the firm}: risk budgets, the factor cap, drawdown stops.
\omcode{code/firm/multistrat/firm_multistrat.py}{59}{95}{Risk budgets with a factor cap, and the firm with drawdown limits.}\label{lst:s1:running-a-multi-strategy-book:firm}
\item \textbf{False stops and netting}.
\omcode{code/firm/multistrat/firm_multistrat.py}{98}{110}{The chance of a false stop, and netting.}\label{lst:s1:running-a-multi-strategy-book:stops}
\item \textbf{Run} \texttt{correlations()}, \texttt{overlap(cap)}, \texttt{limits(limit)}, \texttt{net\_savings()} and \texttt{fig\_multistrat.py}.
\end{tutsteps}
\textbf{What to change next.} Replace the cliff with Grossman and Zhou's proportional rule; estimate the shared factors from P\&L alone and see how
late they are found; plug in the P\&L of this book's own strategies (chapters 2, 5, 19 and 20) as \omterm{def:s1:running-a-multi-strategy-book:multi}{pods}.
\end{tutorial}

\section{Build: the multi-strategy firm}\label{bld:s1:running-a-multi-strategy-book:multistrat}

\begin{build}
\bfield{Purpose} Simulated \omterm{def:s1:running-a-multi-strategy-book:multi}{pods} with shared factors, risk-budget allocation with overlap control, \omterm{def:s1:running-a-multi-strategy-book:limit}{drawdown limits}, false-stop rates and netting.
\bfield{Interface} \texttt{PodConfig(\dots)}, \texttt{simulate\_pods(cfg, rng)}, \texttt{allocate(pnl, load, t, window, factor\_cap)}, \texttt{run\_firm(S,
rebalance, factor\_cap, limit)}, \texttt{stop\_rate(sr, vol, limit, years, n, rng)}, \texttt{netting(trades)}.
\bfield{Rules} Allocations from past data; loadings from positions; stopped \omterm{def:s1:running-a-multi-strategy-book:multi}{pods} restart yearly.
\bfield{Acceptance tests} \texttt{code/firm/multistrat/tests/}: shared factors correlate \omterm{def:s1:running-a-multi-strategy-book:multi}{pods} and crash them together, the cap lowers loaded \omterm{def:s1:running-a-multi-strategy-book:multi}{pods}'
weights, netting by hand, false stops falling with the limit.
\bfield{Stretch} Proportional drawdown control; capital reallocation after stops; the book's own strategies as \omterm{def:s1:running-a-multi-strategy-book:multi}{pods}.
\end{build}

\begin{omsources}
\item A.~E.~Khandani and A.~W.~Lo, ``What happened to the quants in August 2007? Evidence from factors and transactions data'', \emph{Journal of
Financial Markets} 14(1), 2011.
\item S.~J.~Grossman and Z.~Zhou, ``Optimal investment strategies for controlling drawdowns'', \emph{Mathematical Finance} 3(3), 1993.
\item L.~H.~Pedersen, \emph{Efficiently Inefficient: How Smart Money Invests and Market Prices Are Determined}, Princeton University Press, 2015.
\end{omsources}

\section{Exercises}

\begin{exercise}[$\star$]\label{exo:s1:running-a-multi-strategy-book:1}
Ten \omterm{def:s1:running-a-multi-strategy-book:multi}{pods} trade the same instruments with independent trades of equal size. What share of their gross volume does netting save?
\end{exercise}

\begin{exercise}[$\star$]\label{exo:s1:running-a-multi-strategy-book:2}
A \omterm{def:s1:running-a-multi-strategy-book:multi}{pod} loads 0.4 on a factor that falls six \omterm{def:s1:running-a-multi-strategy-book:multi}{pod} volatilities in a crash. What does the \omterm{def:s1:running-a-multi-strategy-book:multi}{pod} lose, in volatilities and in percent at 10\% volatility?
\end{exercise}

\begin{exercise}[$\star$]\label{exo:s1:running-a-multi-strategy-book:3}
Why do three \omterm{def:s1:running-a-multi-strategy-book:multi}{pods} with a correlation of 0.22 lose together in a crash?
\end{exercise}

\begin{exercise}[$\star\star$]\label{exo:s1:running-a-multi-strategy-book:4}
A 5\% \omterm{def:s1:running-a-multi-strategy-book:limit}{drawdown limit} stops 91.6\% of \omterm{def:s1:running-a-multi-strategy-book:multi}{pods} at 10\% volatility within a year. What share does it stop for \omterm{def:s1:running-a-multi-strategy-book:multi}{pods} run at 5\% volatility, and why?
\end{exercise}

\begin{exercise}[$\star\star$]\label{exo:s1:running-a-multi-strategy-book:5}
Why does a tight limit catch the dead \omterm{def:s1:running-a-multi-strategy-book:multi}{pod} fastest, and at what cost?
\end{exercise}

\begin{exercise}[$\star\star$]\label{exo:s1:running-a-multi-strategy-book:6}
Explain Grossman and Zhou's rule and how a two-step \omterm{def:s1:running-a-multi-strategy-book:multi}{pod} limit approximates it.
\end{exercise}

\begin{exercise}[$\star\star\star$]\label{exo:s1:running-a-multi-strategy-book:7}
\emph{Coding.} Run \texttt{overlap(0.25)} with a stricter cap. What happens to the crashes and the Sharpe ratio, and what would you worry about in a real
firm?
\end{exercise}

\begin{exercise}[$\star\star\star$]\label{exo:s1:running-a-multi-strategy-book:8}
\emph{Find the flaw.} ``Our ten \omterm{def:s1:running-a-multi-strategy-book:multi}{pods} have an average pairwise correlation of 0.05 over three years, so the firm's risk is diversified.''
\end{exercise}

\section{Problem: Five at Once}

\begin{problem}[{Weekend problem --- a firm of pods}]\label{pb:s1:running-a-multi-strategy-book:1}
The chapter's synthetic firm and the public record.

\medskip\noindent\textbf{Part I --- The firm.}
\begin{enumerate}
\item Define a \omterm{def:s1:running-a-multi-strategy-book:multi}{multi-strategy book}, a \omterm{def:s1:running-a-multi-strategy-book:multi}{pod} and \omterm{def:s1:running-a-multi-strategy-book:correlation}{strategy correlation}.
\item Describe the synthetic firm: \omterm{def:s1:running-a-multi-strategy-book:multi}{pods}, shared factors, crashes and the dead \omterm{def:s1:running-a-multi-strategy-book:multi}{pod}.
\item How does the firm allocate risk?
\item What did Khandani and Lo find about August 2007?
\end{enumerate}

\medskip\noindent\textbf{Part II --- Shared risk.}
\begin{enumerate}[resume]
\item Give the correlations measured over twenty years and the \omterm{def:s1:running-a-multi-strategy-book:multi}{pods}' crash losses.
\item How does the firm detect the overlap?
\item Give the firm's Sharpe ratio, largest drawdown and crash losses with and without the cap.
\item How much does netting save?
\end{enumerate}

\medskip\noindent\textbf{Part III --- Stops.}
\begin{enumerate}[resume]
\item Define a \omterm{def:s1:running-a-multi-strategy-book:limit}{drawdown limit} and its two errors.
\item Give the false-stop rates on simulated paths and in the firm.
\item How long does the dead \omterm{def:s1:running-a-multi-strategy-book:multi}{pod} run under each limit?
\item Why is the firm's best limit not a robust optimum?
\end{enumerate}

\medskip\noindent\textbf{Part IV --- The verdict.}
\begin{enumerate}[resume]
\item State the \emph{named result}: the firm's drawdown with and without overlap detection, and the false-stop rate of a 5 per cent \omterm{def:s1:running-a-multi-strategy-book:limit}{drawdown limit}.
\item What did Grossman and Zhou show?
\item In what units should limits be set?
\item Why do shared crashes raise false stops at wide limits?
\item What data must every \omterm{def:s1:running-a-multi-strategy-book:multi}{pod} report to the centre?
\item How does this chapter relate to Book~7, chapter~28?
\item Which \omterm{def:s1:running-a-multi-strategy-book:multi}{pods} would you cut first after a crash, and why not all?
\item In one sentence: what does the centre of a multi-strategy firm add?
\end{enumerate}
\end{problem}

\section{Interview questions}

\begin{interviewq}[$\star$ \iqroles{risk}]\label{iq:s1:running-a-multi-strategy-book:1}
How would you set a \omterm{def:s1:running-a-multi-strategy-book:limit}{drawdown limit} for a new \omterm{def:s1:running-a-multi-strategy-book:multi}{pod}?
\end{interviewq}

\begin{interviewq}[$\star\star$ \iqroles{risk}]\label{iq:s1:running-a-multi-strategy-book:2}
Two \omterm{def:s1:running-a-multi-strategy-book:multi}{pods} in different asset classes lost 8\% each in the same week. What do you look for?
\end{interviewq}

\begin{interviewq}[$\star\star$ \iqroles{researcher}]\label{iq:s1:running-a-multi-strategy-book:3}
How would you estimate whether a strategy's edge has died, and how long would it take?
\end{interviewq}

\begin{interviewq}[$\star\star$ \iqroles{developer}]\label{iq:s1:running-a-multi-strategy-book:4}
Design the netting of trades across \omterm{def:s1:running-a-multi-strategy-book:multi}{pods} trading the same instruments.
\end{interviewq}

\begin{interviewq}[$\star\star$ \iqroles{trader}]\label{iq:s1:running-a-multi-strategy-book:5}
Your \omterm{def:s1:running-a-multi-strategy-book:multi}{pod} is 4\% below its peak and the limit is 5\%. What do you do, and what should the firm want you to do?
\end{interviewq}

\begin{interviewq}[$\star\star\star$ \iqroles{researcher}]\label{iq:s1:running-a-multi-strategy-book:6}
A \omterm{def:s1:running-a-multi-strategy-book:multi}{pod}'s cumulative P\&L is a Brownian motion with drift $\mu$ and volatility $\sigma$. Using the fact that for $\mu = 0$ the drawdown from the running
maximum at time $T$ is distributed as $|W_T|$, give the probability that the drawdown exceeds $L$ at a given time $T$, and explain why the probability that
it ever exceeds $L$ within $T$ is larger.
\end{interviewq}
